\section{Cohomological detection and the min--max input}\label{sec:minmax}

This section establishes the topological and variational statements used
below, independently of the geometric family constructed in
Section~\ref{sec:family}. Its output will be a stationary integral
varifold whose mass lies between a fixed threshold above the equator
and the largest mass in the input family. We first identify the
topological condition that forces this lower threshold, then match it
to the published min--max theorem. Throughout, coefficients for chains and
cohomology are $\mathbb Z_2$, and
\[
 p=m+3.
\]
Write $\mathbf M$ for mass and $\mathcal F$ for the flat distance of
chains in $S^{m+1}$. Let $\mathcal C(S^{m+1})$ be the space of sets of
finite perimeter, identified modulo sets of round volume zero, with metric
\[
 d(E,F)=\Vol(E\triangle F).
\]
The mod-two chain carried by $E$ is denoted by $[E]$. The
finite-perimeter representation theorem gives
\begin{equation}\label{eq:minmax-perimeter}
 \mathbf M(\partial[E])=\operatorname{Per}(E),\qquad
 \mathcal F(\partial[E]-\partial[F])\le d(E,F).
\end{equation}
Indeed, the reduced boundary has multiplicity one, so reduction modulo
two does not change its mass \cite[\S2.2]{BDFSphere}.
Finite-mass top-dimensional mod-two chains in the sphere are represented
by measurable sets, with addition given by symmetric difference
\cite[\S7, pp.~317--318]{Federer1978}; see also
\cite[\S\S2.1--2.3 and Appendix]{White2009}.

We record why finite boundary mass gives finite perimeter in the precise
case needed for fillings. Put \(n=m+1\), let \(A=[E]\) be such a chain,
and suppose \(B=\partial A\) has finite mass. For a smooth tangent field
\(X\) on \(S^n\), let \(\phi_t\) be its flow and
\(H(s,x)=\phi_s(x)\). The flat-chain homotopy formula, valid also with
mod-two coefficients \cite[\S6.2, Lemma~6.4]{WhiteGMT}, gives
\[
 (\phi_t)_\#A-A
 =\partial H_\#([0,t]\times A)+H_\#([0,t]\times B)
 =H_\#([0,t]\times B).
\]
The first swept chain vanishes because its dimension is \(n+1\),
whereas its image lies in \(S^n\). On the rectifiable boundary chain,
the tangential \(n\)-Jacobian of \(H\) is at most
\(\|X\|_\infty(1+O_X(t))\): its time derivative is \(X\), and its
\(n-1\) spatial derivatives are those of a flow tending smoothly to the
identity. The area formula for the swept chain therefore gives
\[
 \Vol(\phi_t(E)\triangle E)
 \le t(1+O_X(t))\|X\|_\infty\mathbf M(B).
\]
Since \(\Vol(\phi_t(E))=\int_E J\phi_t\), division by \(t>0\) and
passage to zero show
\(\left|\int_E\operatorname{div}X\right|
 \le\|X\|_\infty\mathbf M(B)\).
The variational definition of perimeter gives
\(\operatorname{Per}(E)\le\mathbf M(B)<\infty\).
The reduced-boundary representation then proves equality of these
masses as in \eqref{eq:minmax-perimeter}. Thus this argument applies to
arbitrary finite-mass fillings before any boundary regularity is known.

\subsection{The filling cover and projective detection}

Denote by $\mathcal Z_m^0(S^{m+1};\mathbb Z_2)$ the flat path component
of zero among finite-mass integral $m$-cycles.

\begin{lemma}[The filling cover]\label{lem:minmax-cover}
The boundary map
\[
 \pi:\mathcal C(S^{m+1})\longrightarrow
 \mathcal Z_m^0(S^{m+1};\mathbb Z_2),\qquad
 \pi(E)=\partial[E],
\]
is a two-sheeted covering map for the indicated topologies. Its deck
involution is $E\mapsto E^c=S^{m+1}\setminus E$.
\end{lemma}

\begin{proof}
We use the small-filling theorem in the following precise form: there
are constants $\nu>0$ and $c\ge1$, depending on the sphere, such that
if $\mathcal F(T-T')<\nu$, then there is an $(m+1)$-chain $R$ in
the sphere with
\begin{equation}\label{eq:minmax-small-filling}
 \partial R=T-T',\qquad
 \mathbf M(R)\le c\mathcal F(T-T').
\end{equation}
This is the isoperimetric choice described in
\cite[\S3.1, p.~11]{MarquesNeves2017}; all locators for this source
refer to arXiv:1311.6501v2. Applied to a sufficiently fine partition of
a flat-continuous path from zero to $T$, these fillings sum to a
top-dimensional chain with boundary $T$. The preceding representation
therefore writes $T=\partial[E]$ with $E$ of finite perimeter.
Conversely, intersecting any such $E$ with the increasing spherical
caps $\{y_{m+2}<t\}$, $-1\le t\le1$, gives sets of finite perimeter
continuous in volume, from the empty set to $E$. Thus every
$\partial[E]$ belongs to this path component.

If $\partial[E]=\partial[F]$, then
$\partial[E\triangle F]=0$. By \eqref{eq:minmax-perimeter}, the
characteristic function of $E\triangle F$ has zero distributional
gradient. Connectedness of the sphere makes it constant almost
everywhere. Hence $F=E$ or $F=E^c$.

For the covering topology, fix $T=\partial[E]$ and choose $\delta>0$
so that $2\delta<\nu$ and $4c\delta<\sigma_{m+1}$. For
$T'$ in the flat $\delta$-ball about $T$, choose the small filling
$R(T,T')$ from \eqref{eq:minmax-small-filling}. This choice is unique:
two such choices differ by zero or the whole sphere, and their total
mass is smaller than $\sigma_{m+1}$. If $T'$ and $T''$ are in the
ball, the same argument applied to
\[
 R(T,T')+R(T,T'')+R(T',T'')
\]
shows that this cycle is zero: its mass is at most $4c\delta$.
Consequently the filling branch determined by
\[
 [E_{T'}]=[E]+R(T,T')
\]
satisfies
$d(E_{T'},E_{T''})\le c\mathcal F(T'-T'')$.
It and its complementary branch are continuous and exhaust the
preimage of the ball. They are disjoint open branches in that
preimage, distinguished by volume distance from $E$ less than, or
greater than, $\sigma_{m+1}/2$. Together with
\eqref{eq:minmax-perimeter}, this proves the assertion.
\end{proof}

Let
\[
 \lambda\in H^1\bigl(\mathcal Z_m^0(S^{m+1};\mathbb Z_2);
                     \mathbb Z_2\bigr)
\]
be the class of this cover: its value on a loop is one precisely when
a lifted filling path ends at the complement of its initial filling.
This description also identifies $\lambda$ with the class in the usual
Almgren--Pitts sweepout definition. Indeed, partition a loop into small
flat increments and lift it to $E_t$. The sum of its small fillings is
$[E_1]+[E_0]$, which is zero or $[S^{m+1}]$ according to the
monodromy. This is exactly Almgren's homology class of the loop
\cite[\S\S3.1--3.3 and the beginning of \S4]{MarquesNeves2017}.

\begin{definition}[Detection]\label{def:detecting}
A flat-continuous map
$\Phi:X\to\mathcal Z_m^0(S^{m+1};\mathbb Z_2)$ from a finite CW
complex is \emph{detecting} if
\[
 (\Phi^*\lambda)^p\ne0\quad\hbox{in }H^p(X;\mathbb Z_2).
\]
Detection is preserved by flat homotopies. We say that $\Phi$ has
\emph{no concentration of mass} if
\begin{equation}\label{eq:minmax-no-concentration}
 \lim_{\rho\downarrow0}\ 
 \sup_{x\in X,\,z\in S^{m+1}}
 \|\Phi(x)\|(B_\rho(z))=0.
\end{equation}
\end{definition}

The following criterion records exactly the topology needed for the
family in Section~\ref{sec:family}.

\begin{lemma}[Projective criterion]\label{lem:minmax-projective}
Set $Q=\overline B^{m+2}\times[-1,1]$. Suppose
$q\mapsto E_q\in\mathcal C(S^{m+1})$ is continuous in volume,
with
\[
 E_{-q}=E_q^c\quad(q\in\partial Q),\qquad
 E_{(0,-1)}=\varnothing,\quad E_{(0,1)}=S^{m+1}.
\]
Then $q\mapsto\partial[E_q]$ descends to a detecting map on
$Q/(q\sim-q\text{ on }\partial Q)\cong\mathbb{RP}^{p}$.
\end{lemma}

\begin{proof}
Radial projection identifies the centrally symmetric convex body $Q$
with the closed $p$-ball, respecting antipodes. Its boundary-antipodal
quotient is therefore $\mathbb{RP}^{p}$. Complementary fillings have
the same mod-two boundary, so \eqref{eq:minmax-perimeter} and the
quotient property give a flat-continuous descended map $\Phi$.
The fillings themselves specify local branches of its pulled-back
cover, not a single filling on the quotient. At a boundary seam one
uses $E_q$ in one collar and the complementary choice in the collar
of $-q$; these choices agree along the identified boundary and are
continuous in volume.

The segment $b\mapsto(0,b)$ becomes a closed loop in this quotient.
Its displayed lift starts at the empty set and ends at the whole
sphere. Hence $\Phi^*\lambda$ is the nonzero degree-one class
$\alpha$ of $\mathbb{RP}^{p}$. The projective cohomology ring is
\[
 H^*(\mathbb{RP}^{p};\mathbb Z_2)
   =\mathbb Z_2[\alpha]/(\alpha^{p+1}),\qquad |\alpha|=1
\]
\cite[Theorem~3.19, p.~220]{Hatcher2002}. In particular
$(\Phi^*\lambda)^p=\alpha^p\ne0$.
\end{proof}

\subsection{A strict perimeter gap}

\begin{lemma}[Detection gap]\label{lem:detection-gap}
There is a constant $\gamma_m>\sigma_m$, depending only on $m$,
such that every detecting map satisfies
\[
 \sup_{x\in X}\mathbf M(\Phi(x))\ge\gamma_m.
\]
\end{lemma}

\begin{proof}
First consider the constrained class of finite-perimeter sets
\[
 \mathcal A_m=\left\{E\subset S^{m+1}:
   \Vol(E)=\frac{\sigma_{m+1}}2,\quad
   \int_E y\,d\Vol(y)=0\right\}.
\]
It is nonempty: the union of two antipodal caps, with their common
radius chosen to give half the spherical volume, belongs to it.
Define
\begin{equation}\label{eq:minmax-gamma}
 \gamma_m=\inf_{E\in\mathcal A_m}\operatorname{Per}(E).
\end{equation}
This number is finite. A minimizing sequence has bounded perimeter;
compactness for sets of finite perimeter on the compact sphere yields
a subsequence converging in $L^1$ to a set $E$ of finite perimeter.
The volume and all centroid coordinates pass to the limit, while
perimeter is lower semicontinuous. Thus $E\in\mathcal A_m$ realizes
\eqref{eq:minmax-gamma}. These are the usual BV compactness and
lower-semicontinuity statements in local charts
\cite[Chapter~2, Theorem~2.6, pp.~54--55]{SimonGMT}; a finite chart cover gives
the compact-sphere assertion used here.

Spherical isoperimetry says that a half-volume set has perimeter at
least $\sigma_m$, with equality only for a hemisphere modulo null
sets \cite[\S2.3, equation~(2.2), p.~5]{BDFSphere}. The centroid of
the hemisphere $\{y\cdot a>0\}$, $|a|=1$, is
\[
 \int_{y\cdot a>0} y\,d\Vol(y)=\frac{\sigma_m}{m+1}a\ne0.
\]
The minimizer in $\mathcal A_m$ therefore cannot attain
$\sigma_m$, proving $\gamma_m>\sigma_m$.

Now pull the filling cover back by $\Phi$. On this cover the function
\begin{equation}\label{eq:minmax-odd-map}
 (x,E)\longmapsto
 \left(\Vol(E)-\frac{\sigma_{m+1}}2,
       \int_E y\,d\Vol(y)\right)\in\mathbb R^{m+3}
\end{equation}
is continuous and changes sign under $E\mapsto E^c$. If it never
vanished, normalization would give an equivariant map to $S^{p-1}$
and hence a quotient map $f:X\to\mathbb{RP}^{p-1}$. Fiberwise this
identifies the filling cover with the pullback of
$S^{p-1}\to\mathbb{RP}^{p-1}$: each of its two points is sent to
one of the two antipodal unit vectors. It follows that
$\Phi^*\lambda=f^*\alpha$, so
$(\Phi^*\lambda)^p=f^*(\alpha^p)=0$, contrary to detection.
At a zero of \eqref{eq:minmax-odd-map} the filling belongs to
$\mathcal A_m$, and its cycle has mass
$\operatorname{Per}(E)\ge\gamma_m$ by
\eqref{eq:minmax-perimeter}.
\end{proof}

\subsection{Cubical parameters and the precise min--max consequence}

\begin{lemma}[Cubical replacement]\label{lem:minmax-cubical}
A detecting map on $\mathbb{RP}^{p}$ can be pulled back to a
detecting map on a finite cubical complex, preserving its mass
supremum and its no-concentration property.
\end{lemma}

\begin{proof}
Embed $\mathbb{RP}^{p}$ in the Euclidean space of real symmetric
$(p+1)\times(p+1)$ matrices by $[u]\mapsto uu^t$, $|u|=1$.
Each such matrix has a simple top eigenvalue $1$, separated from
the remaining eigenvalues $0$. On a neighborhood $U$ of this
compact image the top eigenvalue remains simple, and its eigenline
gives a continuous retraction to $\mathbb{RP}^{p}$. Choose a
sufficiently fine finite cubical grid in a larger cube, and let $K$
be the union of the closed cubes meeting the image, together with
their faces. Then $K\subset U$. After an affine rescaling it is a
finite cubical subcomplex of a subdivided unit cube.

Let $i:\mathbb{RP}^{p}\to K$ be the inclusion and
$r:K\to\mathbb{RP}^{p}$ the restricted retraction. Since
$r\circ i=\mathrm{id}$, the map $r^*$ is injective on cohomology.
Thus $\Phi\circ r$ detects whenever $\Phi$ does. Surjectivity of
$r$ makes the two cycle images identical, proving both remaining
assertions. The dimension of $K$ may exceed $p$. Only its finiteness
and its detecting cohomology class are needed here.
\end{proof}

\begin{theorem}[Min--max consequence]\label{thm:minmax}
Let $m\ge2$, and let
$\Phi:X\to\mathcal Z_m^0(S^{m+1};\mathbb Z_2)$ be a detecting
map on a finite cubical complex, with no concentration of mass and
\[
 D_0=\sup_{x\in X}\mathbf M(\Phi(x))<\infty.
\]
There is a stationary integral $m$-varifold $V$ in $S^{m+1}$ such
that
\begin{equation}\label{eq:minmax-output}
 \gamma_m\le\|V\|(S^{m+1})\le D_0.
\end{equation}
If $2\le m\le6$, $V$ can be chosen as a finite sum of smooth
closed embedded minimal hypersurfaces with positive integer
multiplicities. The same conclusions hold for a family on
$\mathbb{RP}^{p}$.
\end{theorem}

\begin{proof}
We specify both the imported results and the width comparison. For
a discrete map $\phi$ on the vertices of a cubical subdivision,
its \emph{fineness} is
\[
 \mathbf f(\phi)=\max\{
   \mathbf M(\phi(v)-\phi(w)):v,w\text{ are adjacent vertices}\}.
\]
An $(X,\mathbf M)$-homotopy sequence consists of uniformly
mass-bounded discrete maps whose successive maps are joined, after
subdivision, by discrete homotopies of fineness tending to zero.
Two sequences represent the same class if their corresponding maps
are joined by such homotopies with fineness tending to zero. These
are the conventions of
\cite[\S\S2.2--2.6, pp.~6--7]{MarquesNeves2017}.

The discretization theorem
\cite[Theorem~3.9, p.~13]{MarquesNeves2017}, applied using
\eqref{eq:minmax-no-concentration}, provides such a sequence
$\mathcal S=\{\phi_i\}$ with $\mathbf f(\phi_i)\to0$ and
\[
 \limsup_i\max_v\mathbf M(\phi_i(v))\le D_0.
\]
For any sufficiently fine discrete map, the Almgren extension
$\widehat\phi:X\to\mathcal Z_m^0(S^{m+1};\mathbb Z_2)$ is
continuous in the mass norm and satisfies
\begin{equation}\label{eq:minmax-extension}
 \sup_X\mathbf M(\widehat\phi)
 \le\max_v\mathbf M(\phi(v))+C_X\mathbf f(\phi),
\end{equation}
where $C_X$ is fixed
\cite[Theorem~3.10, p.~14]{MarquesNeves2017}.
Moreover, the extensions of $\phi_i$ are flat homotopic to
$\Phi$ for all sufficiently large $i$
\cite[Corollary~3.12(ii), p.~15]{MarquesNeves2017}.

Let $\Pi$ be the class of $\mathcal S$, and write
\[
 \mathbf L(\Pi)=\inf_{\mathcal S'\in\Pi}\mathbf L(\mathcal S'),
 \qquad
 \mathbf L(\mathcal S')=
   \limsup_i\max_v\mathbf M(\phi_i'(v)),
 \quad \mathcal S'=\{\phi_i'\}.
\]
Certainly $\mathbf L(\Pi)\le D_0$. Fix \emph{one} competitor
$\mathcal S'\in\Pi$. Its homotopies to $\mathcal S$ have
fineness tending to zero. Therefore their Almgren extensions are
flat homotopic for all sufficiently large indices, by
\cite[Corollary~3.12(i), p.~15]{MarquesNeves2017}. Each
$\widehat\phi_i'$ is then detecting. Lemma~\ref{lem:detection-gap}
and \eqref{eq:minmax-extension} give
\[
 \gamma_m\le\sup_X\mathbf M(\widehat\phi_i')
 \le\max_v\mathbf M(\phi_i'(v))+C_X\mathbf f(\phi_i').
\]
Taking the upper limit proves
$\gamma_m\le\mathbf L(\mathcal S')$. Since the competitor was
arbitrary,
\[
 0<\gamma_m\le\mathbf L(\Pi)\le D_0.
\]
The index beyond which the extensions detect may depend on
$\mathcal S'$; no threshold uniform over all competitors is used.
This is the continuous/discrete detection compatibility also
recorded in \cite[Lemma~4.6, p.~16]{MarquesNeves2017}.

The Almgren--Pitts existence theorem in
\cite[Theorem~2.14, p.~10]{MarquesNeves2017} now gives a stationary
\emph{integral} varifold at mass $\mathbf L(\Pi)$, almost minimizing
in annuli. That theorem has no dimension restriction. Its
low-dimensional regularity consequence is
\cite[Theorem~2.11, pp.~8--9]{MarquesNeves2017}, which gives the
stated smooth embedded support for $2\le m\le6$. These results
are the mod-two form of the Almgren--Pitts theory; the integrality
input is \cite[Theorem~3.13]{Pitts1981}, and the regularity
extension is \cite[Theorem~4]{SchoenSimon1981}. This proves
\eqref{eq:minmax-output}. For projective parameters, first apply
Lemma~\ref{lem:minmax-cubical}.
\end{proof}

\begin{remark}\label{rem:minmax-scope}
For $m=2,3$, the smooth output in Theorem~\ref{thm:minmax} has a
closed associated mod-two chain: each closed component contributes
its integer multiplicity modulo two. This supplies the cyclicity
needed in Section~\ref{sec:cones}. For $m\ge4$ we use only
stationarity and integrality of the output; regularity will follow
from the density argument there. No min--max index estimate or
multiplicity-one theorem is used. In particular, increasing the
dimension of the cubical parameter complex imposes no new index
requirement. The local index deformation in
Section~\ref{sec:family} is performed on the projective family
before this cubical replacement.
\end{remark}
